\section{TYPE-003: Unconstrained Generic Parameters in toAgentTool}
\label{sec:type-003}
\subsection*{Metadata}
\begin{tabular}{ll}
\textbf{Issue ID} & TYPE-003 \\
\textbf{Severity} & MEDIUM \\
\textbf{Category} & Type Safety \\
\textbf{Branch} & \branchref{fix/TYPE-003-constrain-generic-in-toAgentTool}
\textbf{Changes} & \branchdiff{fix/TYPE-003-constrain-generic-in-toAgentTool} \\
\end{tabular}

\subsection*{Executive Summary}
The \code{toAgentTool} function at \srcfile{src/tool-utils.ts:7} returns \code{AgentTool<any>}, discarding the TypeBox schema type. The caller loses all type information about the tool's parameters. When the tool is executed, the \code{params} argument is typed as \code{any}, allowing invalid parameters to be passed without compile-time errors.

\subsection*{Root Cause Analysis}
The function signature uses \code{any} instead of propagating the schema type:

\begin{lstlisting}[language=TypeScript]
// BEFORE — unconstrained generic and any params
export function toAgentTool(def: GCToolDefinition): AgentTool<any> {
    const schema = buildTypeboxSchema(def.inputSchema);
    return {
        name: def.name,
        parameters: schema,
        execute: async (
            _toolCallId: string,
            params: any,  // No type checking
            signal?: AbortSignal,
        ) => {
            const result = await def.handler(params, signal);
            // ...
        },
    };
}
\end{lstlisting}

\subsection*{Fix Implementation}
The return type is tightened and params typed as \code{Record<string, unknown>}:

\begin{lstlisting}[language=TypeScript]
// AFTER — constrained return type
export function toAgentTool(def: GCToolDefinition): AgentTool<ReturnType<typeof buildTypeboxSchema>> {
    const schema = buildTypeboxSchema(def.inputSchema);
    return {
        name: def.name,
        parameters: schema,
        execute: async (
            _toolCallId: string,
            params: Record<string, unknown>,  // Now typed
            signal?: AbortSignal,
        ) => {
            const result = await def.handler(params, signal);
            // ...
        },
    };
}
\end{lstlisting}

\subsection*{Affected Files}
\begin{itemize}
    \item \srcfile{src/tool-utils.ts} — Lines 4, 14
    \item \srcfile{src/\_\_tests\_\_/type-003.test.ts}
\end{itemize}

\subsection*{Verification}
TypeScript compilation verifies the return type is now constrained.
